\documentclass[12pt]{article}
\usepackage[utf8]{inputenc}
\usepackage[T1]{fontenc}
\usepackage{hyperref}
\title{Uhami Language Resources for K-12 Curriculum Development}
\author{Compiled by Claude Sonnet 5 (claude-sonnet-5) for Andr\'as Kornai}
\date{}
\begin{document}
\maketitle

\section{Introduction}
Uhami (ISO 639-3 uha) is spoken in Ondo State, Nigeria, with about 14,000 speakers (a 2000 count cited by Wikipedia). Its classification is disputed: Edoid in some sources, Akokoid in a newer Wikipedia version. It is sometimes considered the same language as Iyayu. Almost nothing educational was found.

\section{Primary Educational Resources}
None found. The language is not known to be taught in schools.

\section{Linguistic and Reference Materials}
\url{https://en.wikipedia.org/wiki/Uhami_language}

Elugbe's (1989) Edoid classification mentions Uhami among other languages (Glottolog bibliographic record): \url{https://glottolog.org/resource/reference/id/91279}

\url{https://en.wikipedia.org/wiki/Edoid_languages}

\section{Religious and Community Resources}
None found.

\section{Digital and Audio Resources}
None found.

\section{Curriculum Development Context}
No documentation suitable for curriculum use was located. Any work would start with fieldwork and community consultation, and with settling the Uhami/Iyayu and Edoid/Akokoid questions.
\end{document}